% TODO make hw stylesheet
\documentclass[11pt]{scrartcl}
\usepackage[a4paper, total={6in, 8in}]{geometry}
\usepackage[usenames,dvipsnames,svgnames]{xcolor}
\usepackage[shortlabels]{enumitem}
\usepackage[framemethod=TikZ]{mdframed}
\usepackage{amsmath,amssymb,amsthm}
\usepackage[colorlinks]{hyperref}
\usepackage{multicol}
\usepackage{tabularx}

\hypersetup{
	colorlinks=true,
	linkcolor=blue,
	filecolor=magenta,      
	urlcolor=magenta,
	pdftitle={MAT 314 HW}
}

\usepackage{microtype}
\usepackage{mathtools}
\usepackage[headsepline]{scrlayer-scrpage}
\usepackage{thmtools}
\usepackage[textsize=footnotesize]{todonotes}

\mdfdefinestyle{mdbluebox}{roundcorner=10pt,innerbottommargin=9pt,
	linecolor=blue,backgroundcolor=TealBlue!5,}
\declaretheoremstyle[headfont=\sffamily\bfseries\color{MidnightBlue},
mdframed={style=mdbluebox},]{thmbluebox}

\mdfdefinestyle{mdbloobox}{frametitlefont=\bfseries,innerbottommargin=8pt,
	nobreak=true,backgroundcolor=TealBlue!5,linecolor=blue,}
\declaretheoremstyle[headfont=\bfseries\color{MidnightBlue},
mdframed={style=mdbloobox},headpunct={\\[3pt]},postheadspace=0pt,]{thmbloobox}

\mdfdefinestyle{mdredbox}{frametitlefont=\bfseries,innerbottommargin=8pt,
	nobreak=true,backgroundcolor=Salmon!5,linecolor=RawSienna,}
\declaretheoremstyle[headfont=\bfseries\color{RawSienna},
mdframed={style=mdredbox},headpunct={\\[3pt]},postheadspace=0pt,]{thmredbox}

\mdfdefinestyle{mdgreenbox}{linecolor=ForestGreen,backgroundcolor=ForestGreen!5,
	linewidth=2pt,rightline=false,leftline=true,topline=false,bottomline=false,}
\declaretheoremstyle[headfont=\bfseries\sffamily\color{ForestGreen!70!black},
mdframed={style=mdgreenbox},headpunct={ --- },]{thmgreenbox}

\mdfdefinestyle{mdorangebox}{linecolor=RedOrange,backgroundcolor=RedOrange!5,
	linewidth=2pt,rightline=false,leftline=true,topline=false,bottomline=false,}
\declaretheoremstyle[headfont=\bfseries\sffamily\color{RedOrange!70!black},
mdframed={style=mdorangebox},headpunct={ --- },]{thmorangebox}

\mdfdefinestyle{mdblackbox}{linecolor=black,backgroundcolor=RedViolet!5!gray!5,
	linewidth=3pt,rightline=false,leftline=true,topline=false,bottomline=false,}
\declaretheoremstyle[mdframed={style=mdblackbox}]{thmblackbox}

\declaretheoremstyle[spaceabove=6pt,spacebelow=6pt]{basehead}

\declaretheorem[style=basehead,name=Problem]{problem}
\declaretheorem[style=basehead,name=Problem,numbered=no]{problem*}
\declaretheorem[style=thmbloobox,name=Setup]{setup} % for config geo
\declaretheorem[style=thmredbox,name=Example,sibling=problem]{example}
\declaretheorem[style=thmredbox,name=$\bigstar$ Problem,sibling=problem]{niceprob}
\declaretheorem[style=thmbluebox,name=Theorem,sibling=problem]{theorem}
\declaretheorem[style=thmbluebox,name=Corollary,sibling=theorem]{corollary}
\declaretheorem[style=thmbluebox,name=Proposition,sibling=theorem]{prop}
\declaretheorem[style=thmbluebox,name=Lemma,numberwithin=problem]{lemma}
\declaretheorem[style=thmbluebox,name=Theorem,numbered=no]{theorem*}
\declaretheorem[style=thmbluebox,name=Lemma,numbered=no]{lemma*}
\declaretheorem[style=thmgreenbox,name=Claim,numberwithin=problem]{claim}
\declaretheorem[style=thmgreenbox,name=Claim,numbered=no]{claim*}
\declaretheorem[style=thmblackbox,name=Remark,numberwithin=problem]{remark}
\declaretheorem[style=thmblackbox,name=Remark,numbered=no]{remark*}
\declaretheorem[style=thmgreenbox,name=Definition,sibling=theorem]{definition}
\declaretheorem[style=thmgreenbox,name=Definition,numbered=no]{definition*}
\declaretheorem[style=thmorangebox,name=Warning,sibling=theorem]{warning}
\declaretheorem[style=thmorangebox,name=Warning,numbered=no]{warning*}
\declaretheorem[style=thmorangebox,name=Note,numbered=no]{note}
\declaretheorem[style=thmblackbox,name=Exercise,sibling=problem]{exercise}
\declaretheorem[style=thmblackbox,name=Exercise,numbered=no]{exercise*}
\declaretheorem[style=thmblackbox,name=Question,sibling=problem]{question}
\declaretheorem[style=thmblackbox,name=Question,numbered=no]{question*}

\addtolength{\textheight}{3.14cm}
\providecommand{\ol}{\overline}
\providecommand{\ul}{\underline}
\providecommand{\eps}{\varepsilon}
\providecommand{\half}{\frac{1}{2}}
\providecommand{\dang}{\measuredangle} %% Directed angle
\providecommand{\CC}{\mathbb C}
\providecommand{\FF}{\mathbb F}
\providecommand{\NN}{\mathbb N}
\providecommand{\QQ}{\mathbb Q}
\providecommand{\RR}{\mathbb R}
\providecommand{\ZZ}{\mathbb Z}
\providecommand{\dg}{^\circ}
\providecommand{\ii}{\item}
\providecommand{\setdiff}{\smallsetminus}
\providecommand{\alert}[1]{{\sffamily\textbf{\textcolor{blue}{#1}}}}
\providecommand{\inv}{^{-1}}
\providecommand{\mb}[1]{\mathbf{#1}}
\providecommand{\dif}{\;d}

\reversemarginpar

\newenvironment{soln}{\begin{proof}[Solution]}{\end{proof}}
\newenvironment{walkthrough}{\noindent \textbf{Walkthrough.}}{\medskip}
\newlist{walk}{enumerate}{3}
\setlist[walk]{label=\bfseries (\alph*)}

\providecommand{\pow}{\operatorname{Pow}}
\providecommand{\vocab}[1]{{\textbf{\textcolor{ForestGreen}{#1}}}}
\providecommand{\lcm}{\operatorname{lcm}}
\providecommand{\abs}[1]{\left|#1\right|}

\usepackage{asymptote}
\begin{asydef}
	defaultpen(fontsize(10pt));
	unitsize(1cm);
	size(8cm); // set a reasonable default
	usepackage("amsmath");
	usepackage("amssymb");
	settings.tex="pdflatex";
	settings.outformat="pdf";
	// Replacement for olympiad+cse5 which is not standard
	import geometry;
	// recalibrate fill and filldraw for conics
	void filldraw(picture pic = currentpicture, conic g, pen fillpen=defaultpen, pen drawpen=defaultpen)
	{ filldraw(pic, (path) g, fillpen, drawpen); }
	void fill(picture pic = currentpicture, conic g, pen p=defaultpen)
	{ filldraw(pic, (path) g, p); }
	// some geometry
	pair foot(pair P, pair A, pair B) { return foot(triangle(A,B,P).VC); }
	pair orthocenter(pair A, pair B, pair C) { return orthocentercenter(A,B,C); }
	pair centroid(pair A, pair B, pair C) { return (A+B+C)/3; }
	// cse5 abbreviations
	path CP(pair P, pair A) { return circle(P, abs(A-P)); }
	path CR(pair P, real r) { return circle(P, r); }
	pair IP(path p, path q) { return intersectionpoints(p,q)[0]; }
	pair OP(path p, path q) { return intersectionpoints(p,q)[1]; }
	path Line(pair A, pair B, real a=0.6, real b=a) { return (a*(A-B)+A)--(b*(B-A)+B); }
	// cse5 more useful functions
	picture CC() {
		picture p=rotate(0)*currentpicture;
		currentpicture.erase();
		return p;
	}
	pair MP(Label s, pair A, pair B = plain.S, pen p = defaultpen) {
		Label L = s;
		L.s = "$"+s.s+"$";
		label(L, A, B, p);
		return A;
	}
	pair Drawing(Label s = "", pair A, pair B = plain.S, pen p = defaultpen) {
		dot(MP(s, A, B, p), p);
		return A;
	}
	path Drawing(path g, pen p = defaultpen, arrowbar ar = None) {
		draw(g, p, ar);
		return g;
	}
\end{asydef}

\usepackage{mathrsfs}

\ihead{\footnotesize MAT 314 HW01}
\ohead{\footnotesize Last compiled \today}

\title{\vspace{-2em}MAT 314 HW01}
\author{\large Aditya Pahuja}
\date{\large Due 02/06}
\begin{document}
\maketitle

\begin{problem*}1.1.
	\begin{enumerate}[(a)]
		\ii Is $GL_n(\CC)$ isomorphic to a subgroup of $GL_{2n}(\RR)$?
		\ii Is $SO_2(\CC)$ a bounded subset of $C^{2\times 2}$?
	\end{enumerate}
\end{problem*}

\begin{soln}
	For (a), consider the map
	\[Z(a + bi) = \begin{pmatrix}
		a&-b\\b&a
	\end{pmatrix},\]
	which is an injective ring homomorphism
	$Z\colon \CC\to\RR^{2\times 2}$:
	that it is an injective homomorphism of the underlying additive groups
	is trivial to check, and it is well-behaved with multiplication because
	\[(a+bi)(c+di)\mapsto \begin{pmatrix}
		a&-b\\b&a
	\end{pmatrix}\begin{pmatrix}
		c&-d\\d&c
	\end{pmatrix}
	= \begin{pmatrix}
		ac-bd&-ad-bc\\ad+bc&ac-bd
	\end{pmatrix} = Z(ac-bd + (ad+bc)i).\]
	
	Now let $A\in GL_n(F)$ for $F = \CC$ or $\RR$
	have entry $A_{i,j}$ in row $i$, column $j$.
	Then, suppose that $\phi$ is a map sending $A = (A_{i,j})\in GL_n(\CC)$
	to $B = (B_{i,j})\in GL_{2n}(\RR)$ with the rule
	\[\begin{pmatrix}
		B_{2i-1,2j-1}&B_{2i-1,2j}\\
		B_{2i,2j-1}&B_{2i,2j}
	\end{pmatrix} = Z(A_{i,j}).\]
	In fact, it's actually a ring homomorphism
	$C^{1\times 1}\to \RR^{2\times 2}$.
	
	That is, if we abbreviate $Z(A_{i,j})$ as $Z_{i,j}$, then
	\[\phi(A) = \begin{pmatrix}
		Z_{1,1} & Z_{1,2} & \cdots & Z_{1,n}\\
		Z_{2,1} & Z_{2,2} & \cdots & Z_{2,n}\\
		\vdots & \vdots & \ddots & \vdots\\
		Z_{n,1} & Z_{n,2} & \cdots & Z_{n,n}\\
	\end{pmatrix}.\]
	From here we can see that for $X,Y\in GL_n(\CC)$,
	\[(\phi(X)\phi(Y))_{i,j} = \sum_{k=1}^n Z(X_{i,k})Z(Y_{k,j})
	= Z\left(\sum_{k=1}^n X_{i,k}Y_{k,j}\right)
	= \phi(XY),
	\]
	so $\phi$ is a homomorphism, and injectivity of $\phi$ follows from
	injectivity of $Z$. This establishes an isomorphism between
	$GL_n(\CC)$ and the subgroup of $GL_{2n}(\RR)$ corresponding to
	the image of $\phi$.
	
	For (b), we observe that
	\[\left\{\begin{pmatrix}
		x & i\sqrt{x^2-1}\\
		-i\sqrt{x^2-1}&x
	\end{pmatrix} : x > 1\right\}\]
	is a subset of $SO_2(\CC)$, since the inverse of each of these matrices
	is its own transpose and their determinants are all $1$,
	and of course since $x$ can be arbitrarily large
	this shows that $SO_2(\CC)$ is unbounded.
\end{soln}

\begin{problem*}3.4.
	The centralizer of $\mb j$ is just the subset of $SU_2$
	whose $\mb i$ and $\mb k$ components are zero.
	Indeed,
	\[a\mb j + b\mb k - c - d\mb i
	= (a + b\mb i + c\mb j + d\mb k)\cdot\mb j
	= \mb j\cdot (a + b\mb i + c\mb j + d\mb k)
	= a\mb j - b\mb k - c + d\mb i\]
	implies that $b = -b = 0$ and $d = -d = 0$.
\end{problem*}

\begin{problem*}5.1.
	Can the image of a one-parameter group in $GL_n$ cross itself?
\end{problem*}

\begin{soln}
	The answer is \framebox{yes}.
	Take $\phi\colon \RR^+\to GL_n$ to be the one-parameter group
	\[\phi(t) = e^{tA},\qquad A = \begin{pmatrix}
		0&-1&0&0&\cdots&0\\
		1&0&0&0&\cdots&0\\
		0&0&1&0&\cdots&0\\
		0&0&0&1&\cdots&0\\
		\vdots&\vdots&\vdots&\vdots&\ddots&\vdots\\
		0&0&0&0&\cdots&1
	\end{pmatrix}.\]
	Then powers of $A$ fix everything except for the top-left
	$2\times 2$ submatrix, which is periodic under exponentiation
	by an integer with period $4$. In particular
	\[e^{tA} = \begin{pmatrix}
		\cos t&-\sin t&0&0&\cdots&0\\
		\sin t&\cos t&0&0&\cdots&0\\
		0&0&e&0&\cdots&0\\
		0&0&0&e&\cdots&0\\
		\vdots&\vdots&\vdots&\vdots&\ddots&\vdots\\
		0&0&0&0&\cdots&e
	\end{pmatrix},\]
	so the image of $\phi$ crosses itself, as for example
	$e^{0A} = e^{2\pi A}$.
\end{soln}

\begin{problem*}6.1.
	Verify the Jacobi identity for the bracket operation
	$[A,B] = AB - BA$.
\end{problem*}

\begin{soln}
	First, observe that
	\[[A, [B,C]] = A(BC-CB) - (BC-CB)A = ABC - ACB - BCA + CBA.\]
	Similarly,
	\begin{align*}
		[B, [C,A]] &= BCA - BAC - CAB + ACB\\
		[C, [A,B]] &= CAB - CBA - ABC + BAC
	\end{align*}
	and we can see easily that every term with a positive coefficient
	pairs off with its negative when adding the three brackets together.
\end{soln}

\begin{problem*}8.4.
	\begin{enumerate}[(a)]
		\ii Write the polynomial equations that define the symplectic group.
		\ii Show that the unitary group $U_n$ can be defined by
		real polynomial equations in the real and imaginary parts
		of the matrix entries.
	\end{enumerate}
\end{problem*}

\begin{soln}
	(a) Writing elements of $GL_{2n}$ as blocks of $n\times n$ matrices
	\[\begin{pmatrix}
		A&B\\C&D
	\end{pmatrix},\]
	we want
	\[\begin{pmatrix}
		A&C\\B&D
	\end{pmatrix}
	\begin{pmatrix}
		0&I\\-I&0
	\end{pmatrix}
	\begin{pmatrix}
		A&B\\C&D
	\end{pmatrix} =
	\begin{pmatrix}
		0&I\\-I&0
	\end{pmatrix},\]
	which gives the equations
	\begin{align*}
		A^tC - C^tA &= 0\\
		B^tD - D^tB &= 0\\
		A^tD - C^tB &= I,
	\end{align*}
	with the fourth equation $B^tC - D^tA = -I$ being redundant
	because it's obtained by taking the negative transpose
	of both sides of the third equation.
	
	(b) Given a matrix in $U_n$ with entries $z_{i,j}$ in row $i$, column $j$,
	we know the relations
	\[\sum_{k=1}^n z_{i,k}\ol{z_{j,k}} =\begin{cases}
		1, & i=j\\
		0, & i\ne j
	\end{cases}\]
	are equivalent to the matrix being unitary, since
	this is exactly enough to characterize the matrix as being the inverse
	of its conjugate transpose.
	Splitting the $z_{i,j}$ into real and imaginary parts
	and using the above equations gives a collection of
	real polynomials that determine $U_n$.
\end{soln}




















\end{document}